% Annexe : la stabilité et l'ARI expliqués simplement (exercice 3.5).

\section{La stabilité et l'ARI, comme dans une classe (exercice 3.5)}

\subsection{Le test : refaire les équipes avec des absents}

Une maîtresse range ses élèves en équipes selon leurs goûts. Un jour, quelques élèves sont absents : elle
refait les équipes avec ceux qui restent. Si les amis qui étaient ensemble se retrouvent encore ensemble,
et ceux qui étaient séparés restent séparés, sa façon de faire les équipes est \textbf{fiable}.

Le notebook fait pareil avec les plaintes :
\begin{center}
\begin{tikzpicture}[
  b/.style={draw=#1, fill=#1!10, rounded corners=5pt, align=center, font=\small, text width=3.9cm, minimum height=1.6cm},
  fl/.style={-{Stealth[length=3mm]}, very thick, color=slate}]
\node[b=slate] (a) {{\Large\faDatabase}\\toutes les plaintes\\\textbf{rangement complet}};
\node[b=accentblue, right=1.1cm of a] (b) {{\Large\faRandom}\\on tire \textbf{90 \%} des plaintes\\au hasard};
\node[b=sectioncolor, right=1.1cm of b] (c) {{\Large\faRedo}\\on refait \textbf{tout} le rangement\\sur ces 90 \%};
\draw[fl] (a) -- (b); \draw[fl] (b) -- (c);
\node[b=darkblue, below=0.7cm of b, text width=8.6cm, minimum height=1.1cm] (d) {{\faBalanceScale}\ \ on compare les deux rangements \textbf{sur les mêmes plaintes} : c'est l'ARI};
\draw[fl] (c.south) |- (d.east);
\draw[fl] (a.south) |- (d.west);
\end{tikzpicture}
\end{center}

\begin{tcbraster}[raster columns=2, raster equal height, raster column skip=8pt, raster every box/.style={unbreakable}]
\begin{conceptbox}[title={\faRedo\ Trois fois}, unbreakable]
\small On recommence trois fois. Chaque tirage repart de la base \textbf{complète} : on reste toujours à
90~\%, on ne descend pas à 80 puis 70~\%. Seules les plaintes tirées changent d'une fois à l'autre.
\end{conceptbox}
\begin{conceptbox}[title={\faDiceThree\ Un hasard fixé}, unbreakable]
\small Les tirages utilisent les graines 1, 2 et 3 : si on relance le notebook, on retombe sur les trois
mêmes tirages, donc sur les mêmes résultats.
\end{conceptbox}
\end{tcbraster}

\subsection{Comment on compte : des paires d'élèves}

On ne compare pas les \textbf{numéros} des équipes : l'équipe \og 1 \fg{} d'un jour peut s'appeler
\og 2 \fg{} le lendemain, ça ne veut rien dire. On regarde des \textbf{paires} : pour chaque paire d'élèves,
les deux rangements sont-ils d'accord pour dire \og ensemble \fg{} ou \og séparés \fg{} ?

\begin{center}
\begin{tikzpicture}[el/.style={circle, draw=#1, fill=#1!15, minimum size=0.8cm, font=\small\bfseries}]
\node[font=\small\bfseries, color=darkblue] at (1.6,2.3) {lundi};
\draw[accentblue, thick, rounded corners=8pt] (-0.1,0.6) rectangle (1.5,1.6);
\draw[accentblue, thick, rounded corners=8pt] (1.7,0.6) rectangle (3.3,1.6);
\node[el=accentblue] at (0.3,1.1) {A}; \node[el=accentblue] at (1.1,1.1) {B};
\node[el=ambre] at (2.1,1.1) {C}; \node[el=ambre] at (2.9,1.1) {D};
\begin{scope}[xshift=5cm]
\node[font=\small\bfseries, color=darkblue] at (1.6,2.3) {mardi};
\draw[accentblue, thick, rounded corners=8pt] (-0.1,0.6) rectangle (1.5,1.6);
\draw[accentblue, thick, rounded corners=8pt] (1.7,0.6) rectangle (2.5,1.6);
\draw[accentblue, thick, rounded corners=8pt] (2.6,0.6) rectangle (3.4,1.6);
\node[el=accentblue] at (0.3,1.1) {A}; \node[el=accentblue] at (1.1,1.1) {B};
\node[el=ambre] at (2.1,1.1) {C}; \node[el=ambre] at (3.0,1.1) {D};
\end{scope}
\node[anchor=north west, font=\small, text width=6.2cm] at (9.6,2.4) {
\begin{tabular}{@{}lll@{}}
\textbf{paire} & \textbf{lundi, mardi} & \\
A et B & ensemble, ensemble & \textcolor{accentblue}{\faCheck} \\
A et C & séparés, séparés & \textcolor{accentblue}{\faCheck} \\
A et D & séparés, séparés & \textcolor{accentblue}{\faCheck} \\
B et C & séparés, séparés & \textcolor{accentblue}{\faCheck} \\
B et D & séparés, séparés & \textcolor{accentblue}{\faCheck} \\
C et D & ensemble, séparés & \textcolor{sectioncolor}{\faTimes} \\
\end{tabular}};
\end{tikzpicture}
\end{center}
{\small\color{slate}\textit{Exemple inventé : 4 élèves, 6 paires ; les deux jours sont d'accord sur 5 paires sur 6.}}

\subsection{Pourquoi \og ajusté \fg{} ?}

Compter les paires d'accord donne l'\textbf{indice de Rand}. Mais même deux rangements faits \textbf{au
hasard} tombent souvent d'accord, surtout sur \og séparés \fg, parce que la plupart des paires sont séparées
dans n'importe quel rangement. L'indice de Rand \textbf{ajusté} (ARI) retire cette part due au hasard :

\begin{center}
\begin{tikzpicture}[x=1.35cm]
\shade[left color=palblue, right color=accentblue!80] (0,0) rectangle (10,0.45);
\draw[slate] (0,0) rectangle (10,0.45);
\foreach \x/\v in {0/0, 5/0{,}5, 10/1} { \draw[slate] (\x,0) -- (\x,-0.15); \node[below, font=\small] at (\x,-0.15) {\v}; }
\node[above, font=\footnotesize, text width=3.4cm, align=center] at (0.9,0.5) {pas mieux que le hasard};
\node[above, font=\footnotesize\bfseries, color=sectioncolor, text width=3.4cm, align=center] at (5,0.5) {seuil du notebook};
\node[above, font=\footnotesize, text width=3.4cm, align=center] at (9.1,0.5) {rangements identiques};
\draw[very thick, dashed, color=sectioncolor] (5,-0.05) -- (5,0.5);
\end{tikzpicture}
\end{center}

\doisvoir{trois valeurs d'ARI, une par tirage, toutes au-dessus de 0,5.}

\begin{piege}
\textbf{Stable ne veut pas dire juste.} Mettre tout le monde dans une seule équipe donne un ARI parfait
à chaque fois, et ne range rien. C'est pourquoi l'exercice 3.3 écarte d'abord les rangements ratés, et ne
regarde la stabilité qu'ensuite.
\end{piege}
